\scanpage{88}
\mspage{\handwritten{64}}

(\textsuperscript{22}) For the notion of a ``lien'' or ``tie'', \deleted{\unsure{see}}which is one of the main ingredients of the non commutative cohomology panoply of Giraud's theory, I refer to his book (Springer Grundlehren 179, 1971). \inserted{A} Picard categor\deleted{ies}y \deleted{\unsure{are categories}}\inserted{is a groupoi\unsure{d}} endowed with an operation $\unsure{\otimes|\oplus}$ \deleted{which makes it \unsure{ressemble|resemble} to an} together with associativity, unity and commutativity data for this operation, which make it reseumble to a commutative group. \inserted{A} ``Champ\deleted{\unsure{s}} de Picard'' or ``Picard stacks'' \deleted{are}\inserted{is} defined accordingly, by relativizing over an arbitrary space or topos (replacing the groupoid by a stack of groupoids over this topos). The necessary generally nonsense on these is devel\unsure{p|o}ped rathe\unsure{r} carefully in an exposé of Deligne in SGA 4 (SGA 4 XVIII 1.4). In this letter to Larry Breen, I am assuming ``known'' the notion of \deleted{\unsure{a}}an $n$-stack (for $n=\deleted{\unsure{2|3}}\unsure{3|2}$ at any rate), and the corresponding notion of (strict) $n$-Picard stack, which should be describable (as was explained in Deligne's notes in the case $n=1$) by an $n$-truncated chain complex in the category of abelian sheaves on $\mathrm{X}$ (viewed mainly as an object of the relevant derived category). The ``strictness'' condition on usual\unsure{,|'} Picard stacks refers to the restriction that the commutativity isomorphism within an object $\mathrm{L}\otimes\mathrm{L}'$ , \deleted{\unsure{wh}} when $\mathrm{L}=\mathrm{L}'$, should reduce to the identity. It is assumed (wi thout further explanation) that the condition carries over in a natural way to Picard $n$-stacks, in such a way as to allow an interpretation of these by truncated objects in a suitable derived category, as hinted above.

(\textsuperscript{23}) When $\mathrm{M}$ is any abelian sheaf on a topos, $\mathrm{X}(\mathrm{M})$ is a certain canonical \inserted{left} resolution of $\mathrm{M}$ by sheaves of $\underline{\mathrm{Z}}$-modules which are ``free'', and more specifically, which are finite direct sums of sheaves of the type $\underline{\mathrm{Z}}^{(\mathrm{T})}$ , where $\mathrm{T}$ is any sheaf of the type $\mathrm{M}^{n}$ (finite product of copies of $\mathrm{M}$ ). This canonical\deleted{ly} construction was introduced \deleted{first} by MacLane, and gained new popularity in the French school of algebraic geometry and homological algebra in th\unsure{ê|e} late sixties, because \deleted{in case M is described by some abelian group scheme over a ground} it gives a very handy way to relate the $\mathrm{Ext}^{i}(\mathrm{M},\mathrm{N})$ invariants (when $\mathrm{N}$ is another abelian sheaf on $\mathrm{X}$) to the \deleted{\unsure{sp}} ``spacial'' cohomology of $\mathrm{M}$ (i.e. of the induced topos $\mathrm{X}_{/\mathrm{M}}$\deleted{\unsure{,}}) with coefficients in $\mathrm{N}$.

(\textsuperscript{24}) Reflecting on the ``right'' version of the provisional Verdier notion of a \deleted{derived}\inserted{triangulated} category (which was supposed to describe adequately the \inserted{relevant} internal structure of the derived categories of abelian categories) is part of my present program for the notes on Pursuing Stacks, and will be the main task in \deleted{\unsure{t}}one of the chapters of volume two.

(\textsuperscript{25}) As was seen in section 9, ``uniqueness'' here has to be understood in
% archive sheet number: 87
